\section{Cross-task requirements and evidence-reuse results}\label{sec:r9-results}
\subsection{Configured service support and complete-task limits}
The complete catalogue contains 192 task/variant/population/service requests per study and allowance. Figure~\ref{fig:r9-support} and Table~\ref{tab:r9-support} retain unsupported requests in the denominator. At the primary allowance, the two semantic implementations support 24 requests in A and 26 in B. Direct renewal supports 30 and 32; exact-anchor matching supports 18 and 20. The equality of the contract and complete procedural implementation holds for all 6,144 frozen plans and their repeated confirmation outputs. The semantic layer therefore supplies a transferable interpretation of evidence, but the results do not establish an advantage exclusive to its representation.

\begin{table}[!htbp]\centering\small
\caption{Primary-allowance configured support. Each denominator retains every task, variant, population and service. Full-task support counts only the original complete requirement.}\label{tab:r9-support}
\begin{tabular}{@{}llll@{}}\toprule
Study & Method & All requests & Full tasks\\\midrule
A & ContractCover & 24/192 & 0/24\\
A & ProceduralCover & 24/192 & 0/24\\
A & DirectRenewal & 30/192 & 0/24\\
A & ExactCard & 18/192 & 0/24\\
B & ContractCover & 26/192 & 1/24\\
B & ProceduralCover & 26/192 & 1/24\\
B & DirectRenewal & 32/192 & 1/24\\
B & ExactCard & 20/192 & 1/24\\
\bottomrule\end{tabular}\end{table}

\begin{figure}[!htbp]\centering\includegraphics[width=\linewidth]{r9/support.pdf}
\caption{Configured support under the full request matrix, with a common confidence allocation, policy objective and records. The primary allowance is 0.15. Contract and complete procedural curves coincide. Direct renewal retains its additional supported services; exact-card matching is a limited mechanism control. These are configuration counts, not a field-request probability.}\label{fig:r9-support}\end{figure}

At allowance 0.15, no complete acquisition task is certified in either study. The only full-task semantic plan is B/unfitted/relaxed association, using FlatCART with upper bound 0.143973. Its margin is 0.006027, so its acceptance is narrow. Complete reporting remains refused in every population/service cell. Most supported requests are partial or single-component services, and several nominal variants compile to the same event. They must not be described as independent successful tasks or as recovery of the full inspection service.

The wider ROI adds two primary supported configurations, yet its larger virtual production cost limits timely reporting. Secondary allowances expose additional configured support, including B's increase from 26 to 30 at allowance 0.20; changing the allowed risk is not an improvement under the original 0.15 requirement. Figure~\ref{fig:r9-matrix} displays every primary request, including the direct-renewal-only services and the unsupported tight/contended region. A stronger direct baseline is retained instead of being constrained to the anchor family.

\begin{figure}[!htbp]\centering\includegraphics[width=\linewidth]{r9/primary_matrix.pdf}
\caption{All primary configurations. Each study has eight population/service rows and 24 task/variant columns. V1--V8 denote full, coarser GSD, more blur, partial, partial relaxed, stricter, late report and single component. Blue indicates semantic support; open orange circles retain additional direct-renewal support. Compiled-event duplicates remain visible but do not increase evidence sample size.}\label{fig:r9-matrix}\end{figure}
\begin{table}[!htbp]\centering\small
\caption{Fixed before-sensing plans and independent confirmation counts; the rows are representative configured services, not independent copies of the source sample. All plans are supplied as CSV.}\label{tab:r9-selected}
\begin{tabular}{@{}llllllll@{}}\toprule
Study & Pop. & Service & Task & Variant & Policy & $U$ & Confirm. $k/n$\\\midrule
A & nom & rel & acqui & partial & OldRAVC & 0.1472 & 110/1024\\
A & nom & rel & acqui & single service & OldRAVC & 0.1375 & 98/1024\\
A & nom & rel & assoc & partial & FlatCART & 0.1494 & 120/1024\\
A & nom & rel & assoc & single service & FlatCART & 0.1128 & 82/1024\\
A & nom & rel & repor & single service & OldRAVC & 0.0069 & 1/1024\\
A & unf & rel & assoc & partial & FlatCART & 0.0733 & 54/1024\\
A & unf & rel & assoc & single service & FlatCART & 0.0568 & 39/1024\\
A & unf & rel & repor & single service & FlatCART & 0.0056 & 1/1024\\
B & nom & rel & acqui & partial & OldRAVC & 0.1391 & 106/1024\\
B & nom & rel & acqui & single service & OldRAVC & 0.1285 & 106/1024\\
B & nom & rel & assoc & partial & FlatCART & 0.1095 & 90/1024\\
B & nom & rel & assoc & single service & OldRAVC & 0.1423 & 123/1024\\
B & nom & rel & repor & single service & OldRAVC & 0.0080 & 1/1024\\
B & unf & rel & assoc & full & FlatCART & 0.1440 & 103/1024\\
B & unf & rel & assoc & partial & FlatCART & 0.0750 & 43/1024\\
B & unf & rel & assoc & single service & FlatCART & 0.0616 & 29/1024\\
B & unf & rel & repor & single service & FlatCART & 0.0069 & 0/1024\\
\bottomrule\end{tabular}\end{table}


\subsection{Optical intervention and causal association}
Actual rasters are produced by raycasting and decoded without truth coordinates. Figure~\ref{fig:r9-rasters} shows the first preselected excluded pilot under the two ROIs. The 48 paired pilot sources verify identical fixed-target GSD: enlarging the ROI does not change range, focal length or incidence at the planned addresses. It can acquire more targets while leaving the sampling-density requirement unsatisfied. Figure~\ref{fig:r9-optical} reports the independent confirmation distributions, separated by their calibrated populations. Study A and B confirmation samples use different seeds; their empirical differences are not a paired intervention estimate.

\begin{figure}[!htbp]\centering\includegraphics[width=.9\linewidth]{r9/paired_rasters.pdf}
\caption{Actual rendered pixels for excluded pilot seed 12000000 under the predeclared ROI intervention. Targets and pixel focal length are fixed. These synthetic diagnostic landmarks illustrate the producer, not aircraft defects or real-camera observations. The same pilot source is used in both panels.}\label{fig:r9-rasters}\end{figure}
\begin{figure}[!htbp]\centering\includegraphics[width=\linewidth]{r9/optical_boundary.pdf}
\caption{Independent confirmation empirical GSD distributions for the nominal and unfitted populations, with a 20-mm/pixel requirement reference. Fixed-target GSD is a geometric score and includes planned targets outside the acquired ROI. The functional ROI invariance is checked on paired pilots; the formal A/B samples are independent.}\label{fig:r9-optical}\end{figure}

Association uses the numerical version available at capture. A later update cannot repair readiness retrospectively. Report delay includes the raster/decoder virtual-service clock; image age subtracts the independently sampled capture time. This distinguishes acquisition, association and delivery instead of measuring three copies of localization accuracy. The calibration and confirmation tables retain all task quantities and every policy/service outcome, including nonfinite-score failures.

\subsection{Certificate stability and transfer boundaries}
All 3,072 stored cards have been recomputed from their executed failure counts with the reserved 4,608-family allocation. Cards identify the population, requested or anchor event, complete producer dependencies and raw calibration lineage. Invariant pooled cards use 4,096 independent inputs; other cards use 2,048. The unused reserve is not redistributed.

Figure~\ref{fig:r9-prefix} retains a passing full-association anchor, a component optical anchor and an unsupported full optical event. Smaller nested prefixes frequently lose certification, and even a full-sample passing association card has limited primary margin. Prefix sizes and alternative confidence levels are diagnostic grids; they do not receive an additional simultaneous guarantee, change F1, or produce renewed confirmation-based selection.
\begin{figure}[!htbp]\centering\includegraphics[width=\linewidth]{r9/certificate_prefix.pdf}
\caption{Actual fixed-event failure counts under nested calibration prefixes and the original reserved-family allocation. The 0.15 line is a requirement, not a fitted target. Displayed prefixes do not create independent samples or an extra guarantee over this sensitivity grid. Full-precision counts and confidence settings are supplied as CSV.}\label{fig:r9-prefix}\end{figure}

Evidence reuse through the ROI change is permitted only for the implemented invariant requested projection. It does not certify new coverage, association or delivery functions by attaching an unchanged optical capsule. Seventy-two excluded provider-change fixtures verify conservative invalidation after capture, motion or geometry provider edits. Thirty-two malformed receipt fixtures produce no incorrect authorization, maintain refusal across reconstructed gates, and admit independent trusted successors. These deterministic software fixtures demonstrate the stated applicability rules, not a probabilistic fault rate or detection of unknown physical drift.

The source-bound receipt observer runs against frozen plans and archived execution records. It binds the trusted source identity, source seal, population, study, policy, compiled event, plan and card. This is executed attachment bookkeeping after the numerical experiment, rather than an added real-time camera pipeline. Its measured cost is reported separately. The source/prior/version enforcement inside the retained localization executor remains a different, actually executed component.

\subsection{Executed update and cache costs}
\begin{table}[!htbp]\centering\small
\caption{Executed computation diagnostics with equally preloaded records and event-identity deduplication. Cold event preparation is the median of three trials; the first batch contains 384 primary queries. Warm time is per shared-cache lookup. These are not developer-effort or embedded worst-case measurements.}\label{tab:r9-cost}
\begin{tabular}{@{}llllll@{}}\toprule
Method & Cards & Record visits & Prep. (s) & First batch (s) & Warm ($\mu$s)\\\midrule
ContractCover & 1,152 & 2,490,368 & 0.179 & 6.408 & 0.091\\
ProceduralCover & 1,152 & 2,490,368 & 0.176 & 6.225 & 0.111\\
DirectRenewal & 2,560 & 5,373,952 & 0.365 & 9.991 & 0.098\\
ExactCard & 1,152 & 2,490,368 & 0.186 & 8.747 & 0.103\\
\bottomrule\end{tabular}\end{table}

With event-identity deduplication available to every method, the anchor methods prepare 1,152 unique cards and the direct method prepares 2,560. Their record visits are 2,490,368 and 5,373,952 respectively, a 53.66\% reduction for anchor reuse. This is a trade-off: less event re-evaluation supports fewer configured requests. The complete procedural method obtains the same reduction. Preparation wall time is measured on equally preloaded records; disk reads, numerical producer execution and human development are outside this comparison.

The first 384-request batch separately includes numeric-card validation and plan search. Every method then receives the same request dictionary cache. Warm lookups require no record rescoring, and their small timing differences are not treated as a scientifically exclusive speed advantage. Known direct-event duplicates are cached as well. JSON arrays are restored to the declared immutable context tuple fields before procedural re-query, so serialization does not artificially reject that baseline. All re-queried plans reproduce F1.

\subsection{Numerical trade-offs, geometry and retained routes}
Figure~\ref{fig:r9-pareto} retains all eight numerical policies, separated by B767, C172 and the unfitted model. The Pareto boundary is descriptive. Fixed2 is kept wherever its work or error trade-off is favorable; a representation claim does not remove a competitive fixed strategy. Complete numerical summaries include timely positive-version success, mean error, mean episode-tail error, 95th percentile of episode maxima, work, publications and native components. Population calibration remains on the frozen mixture, so the model-separated descriptive tables do not imply separately calibrated conditional certificates.
\begin{figure}[!htbp]\centering\includegraphics[width=.96\linewidth]{r9/geometry_pareto.pdf}
\caption{Study B confirmation numerical work--loss means, separately for each geometry at standard and contended services. All eight policies, including Fixed2, remain visible. Camera work is common across policies in each panel and omitted from the horizontal axis only to expose numerical scheduling trade-offs. Full total-work values are supplied. Nondominated means are not confidence bounds.}\label{fig:r9-pareto}\end{figure}

Unsupported occlusion, degenerate range sectors, larger priors and map warps receive no transferred population card. All 512 stress inputs are retained as descriptive evidence. The preserved R8 study additionally contains 36 complete-route inputs and 432 executions. Ordinary routes can favor Fixed2; C172 contention can improve KG1 timely success while worsening its RMSE. These archived route observations are not rerun or renamed as new R9 evidence. They do not certify a route-level image/report event. The original platform photographs and simulation views are included in Section~\ref{sec:design}; the route comparisons and continuous traces are reported in Section~\ref{sec:evaluation}. The supplementary material retains the complete all-method route tables.
